% !TEX program = xelatex
\documentclass[12pt,a4paper]{report}

% ------------------------------------------------------------
% Compile in Overleaf with XeLaTeX.
% The thesis uses Unicode Armenian text and keeps technical terms in English.
% ------------------------------------------------------------

\usepackage[a4paper,left=3cm,right=1.5cm,top=2cm,bottom=2.5cm]{geometry}
\usepackage{fontspec}
\usepackage{polyglossia}
\setmainlanguage{armenian}
\setotherlanguage{english}
\setmainfont{DejaVu Serif}
\newfontfamily\armenianfont{DejaVu Serif}
\newfontfamily\englishfont{DejaVu Serif}
\setmonofont{DejaVu Sans Mono}
\newfontfamily\armenianfonttt{DejaVu Sans Mono}

\usepackage{setspace}
\onehalfspacing
\usepackage{indentfirst}
\setlength{\parindent}{1cm}
\setlength{\parskip}{0pt}

\usepackage{amsmath,amssymb,mathtools}
\usepackage{graphicx}
\usepackage{array,booktabs,longtable,tabularx}
\usepackage{caption}
\usepackage{float}
\usepackage{listings}
\usepackage{xcolor}
\usepackage{hyperref}
\usepackage{titlesec}
\usepackage{tocloft}

\hypersetup{
    colorlinks=true,
    linkcolor=black,
    citecolor=black,
    urlcolor=black
}

\captionsetup[figure]{font=footnotesize,labelfont=bf,justification=centering,labelsep=period}
\captionsetup[table]{font=footnotesize,labelfont=bf,justification=centering,labelsep=period,position=top}

\titleformat{\chapter}[display]
  {\normalfont\bfseries\centering\fontsize{12pt}{18pt}\selectfont}
  {ԳԼՈՒԽ \thechapter}
  {0.5em}
  {\MakeUppercase}
\titlespacing*{\chapter}{0pt}{0pt}{1.5em}

\titleformat{\section}
  {\normalfont\bfseries\fontsize{12pt}{18pt}\selectfont}
  {\thesection}{0.8em}{}
\titleformat{\subsection}
  {\normalfont\bfseries\fontsize{12pt}{18pt}\selectfont}
  {\thesubsection}{0.8em}{}

\renewcommand{\contentsname}{Բովանդակություն}
\renewcommand{\bibname}{ՕԳՏԱԳՈՐԾՎԱԾ ԳՐԱԿԱՆՈՒԹՅԱՆ ՑԱՆԿ}
\renewcommand{\tablename}{Աղյուսակ}
\renewcommand{\figurename}{Նկար}

% Force Western Arabic digits in headings, TOC, pages and enumerate labels
% even under Armenian language settings.
\makeatletter
\newcommand{\UseWesternDigits}{%
  \renewcommand{\thepage}{\number\c@page}%
  \renewcommand{\thechapter}{\number\c@chapter}%
  \renewcommand{\thesection}{\number\c@chapter.\number\c@section}%
  \renewcommand{\thesubsection}{\number\c@chapter.\number\c@section.\number\c@subsection}%
  \renewcommand{\thefigure}{\number\c@chapter.\number\c@figure}%
  \renewcommand{\thetable}{\number\c@chapter.\number\c@table}%
  \renewcommand{\theequation}{\number\c@chapter.\number\c@equation}%
  \renewcommand{\theenumi}{\number\c@enumi}%
  \renewcommand{\theenumii}{\number\c@enumii}%
  \renewcommand{\labelenumi}{\theenumi)}%
  \renewcommand{\labelenumii}{\theenumii)}%
}
\makeatother
\AtBeginDocument{\UseWesternDigits}

\lstdefinestyle{code}{
  basicstyle=\ttfamily\footnotesize,
  breaklines=true,
  frame=single,
  columns=fullflexible,
  showstringspaces=false,
  keepspaces=true
}

\newcommand{\Eng}[1]{\textenglish{#1}}
\newcommand{\TODO}[1]{\textbf{[ԼՐԱՑՆԵԼ՝ #1]}}

\begin{document}

% ------------------------------------------------------------
% TITLE PAGE
% ------------------------------------------------------------
\begin{titlepage}
\begin{center}
{\bfseries Երևանի Պետական Համալսարան}\\[0.3cm]
{\bfseries Ֆիզիկայի Ֆակուլտետ}\\[0.3cm]
% {\bfseries Տվյալների մշակումը ֆիզիկայում և ԱԲ}\\[2.0cm]

{\bfseries ԲԱԿԱԼԱՎՐԻ ԱՎԱՐՏԱԿԱՆ ԱՇԽԱՏԱՆՔ}\\[1.5cm]

{\bfseries Թեմա՝}\\[0.4cm]
{\bfseries \large \Eng{ԱԲ և մեծ լեզվական մոդելների վրա հիմնված նկարների փոխակերպու} համակարգի նախագծում և իրականացում}\\[2.0cm]

\begin{tabular}{>{\bfseries}p{5cm}p{8cm}}
Ուսանող՝ & Սպարտակ Աղաբաբյան \\
Մասնագիտություն՝ & Տվյալների մշակումը ֆիզիկայում և ԱԲ \\
Ղեկավար՝ & Արթուր Կոբելյան \\
\end{tabular}

\vfill
{\bfseries Երևան - 2026}
\end{center}
\end{titlepage}

% ------------------------------------------------------------
% SIGNATURE PAGE
% ------------------------------------------------------------
% \thispagestyle{empty}
% \begin{center}
% {\bfseries ՍՏՈՐԱԳՐՈՒԹՅՈՒՆՆԵՐԻ ԷՋ}
% \end{center}
% \vspace{2cm}

% \begin{tabularx}{\textwidth}{>{\bfseries}p{5cm}X p{4cm}}
% Ուսանող՝ & Սպարտակ Աղաբաբյան & \hrulefill \\
% & & {\centering ստորագրություն\par} \\[1.2cm]
% Ղեկավար՝ & \TODO{ղեկավարի անունը} & \hrulefill \\
% & & {\centering ստորագրություն\par} \\[1.2cm]
% Ամբիոնի վարիչ՝ & \TODO{անունը} & \hrulefill \\
% & & {\centering ստորագրություն\par} \\
% \end{tabularx}

\newpage

% ------------------------------------------------------------
% ABSTRACT
% ------------------------------------------------------------
\thispagestyle{empty}
\begin{center}
{\bfseries ՀԱՄԱՌՈՏԱԳԻՐ}
\end{center}

% \noindent\textbf{Հայերեն անվանում՝} \Eng{AI Image Manipulation} համակարգի նախագծում և իրականացում։\\
% \textbf{Ռուսերեն անվանում՝} Разработка и реализация системы \Eng{AI Image Manipulation}.\\
% \textbf{Անգլերեն անվանում՝} \Eng{Design and Implementation of an AI Image Manipulation System}.

Այս աշխատանքը նվիրված է \Eng{AI Image Manipulation} համակարգի նախագծմանը և իրականացմանը։ Աշխատանքի հիմքում ընկած է այն գաղափարը, որ image editing-ը կարելի է ներկայացնել ոչ միայն որպես պիքսելներիի անմիջական գեներացում, այլ նաև որպես գործողությունների հաջորդականություն, որտեղ մոդելը տեսնում է ներկա նկարը, ստանում է natural language instruction, ապա վերադարձնում է ստրուկտուրացված գործողությունների հաջորդականություն։ Նման մոտեցումը թույլ է տալիս image manipulation խնդիրըը մոտեցնել \Eng{vision-language-action} խնդրին, որտեղ արդյունքը ոչ թե անմիջապես վերջնական պատկերն է, այլ վերահսկելի և վերարտադրելի գործողությունների շարք։

Աշխատանքում նախագծվել է data annotation tool, որը թույլ է տալիս օգտատիրոջը կատարել brush, pencil, eraser, fill, picker, color adjustment, text overlay, blur, clone stamp, scatter brush, pattern brush և forward warp գործողություններ։ Յուրաքանչյուր գործողություն գրանցվում է session log-ում, իսկ նկարի միջանկյալ վիճակները պահպանվում են որպես frames։ Այնուհետև մշակվել է data preparation pipeline, որը raw operations-ը խմբավորում է strokes-ի, սեղմում է անում տրայեկտորիաները \Eng{Ramer--Douglas--Peucker} algorithm-ի միջոցով, բաժանում է տվյալները fixed-horizon action chunks-ի և ձևավորում է \Eng{HuggingFace Dataset}՝ model fine-tuning-ի համար։

Մոդելավորման մասում ներկայացվում է \Eng{Qwen/Qwen3.5-4B} vision-language model-ի վրա հիմնված ուսուցում, որտեղ օգտագործվում է \Eng{QLoRA} supervised fine-tuning և reward-based \Eng{GRPO} training։ Համակարգի inference փուլը կառուցված է observe--predict--execute--reobserve loop-ի վրա․ մոդելը կանխատեսում է action chunk, deterministic canvas executor-ը կիրառում է գործողությունները, ապա ստացված նկարը կրկին տրվում է մոդելին մինչև completion կամ սահմանված iteration limit։

Աշխատանքի ելակետն այն է, որ ժամանակակից image generation model-ները արդեն լավ են լուծում ամբողջական image generation խնդիրը, ուստի ավելի հետաքրքիր է կառուցել controlled editing layer existing նկարի վրա։ Արդյունքը research prototype է, որը միավորում է annotation, data processing, action schema validation, deterministic rendering, supervised training, reward-based optimization և inference pipeline components։ Աշխատանքը ցույց է տալիս, որ image editing task-ը կարելի է ձևակերպել որպես controllable tool-based action generation problem, ինչը ավելի թափանցիկ և վերլուծելի մոտեցում է՝ համեմատած ամբողջական image-to-image black-box generation-ի հետ։

\newpage

% Page numbering starts from the table of contents; according to the standard it is counted as page 4.
\pagenumbering{arabic}
\UseWesternDigits
\setcounter{page}{4}
\tableofcontents
\newpage

% ------------------------------------------------------------
\chapter*{ՆԵՐԱԾՈՒԹՅՈՒՆ}
\addcontentsline{toc}{chapter}{ՆԵՐԱԾՈՒԹՅՈՒՆ}

Այսօր image generation model-ները արդեն հասել են այնպիսի որակի, որ այս աշխատանքի շրջանակում 0-ից general-purpose image generation model կառուցելը նպատակահարմար չէ։ Խնդիրը ոչ թե նոր text-to-image generator ստեղծելն է, այլ ուսումնասիրելը, թե ինչպես կարելի է existing նկարի վրա կատարել վերահսկվող, լոկալ և մեկնաբանելի փոփոխություններ՝ առանց ամբողջ նկարը անտեղի վերագեներացնելու։

Դասական image generation կամ image-to-image editing մոտեցումների հիմնական սահմանափակումներից մեկն այն է, որ ինստրուկցիա ստանալիս model-ը հաճախ փոփոխում է ոչ միայն խնդրված հատվածը, այլև նկարի այլ մասերը։ Օրինակ, եթե օգտատերը խնդրում է blur անել logo-ն, նկարել մազեր կամ փոխել միայն փոքր տեղամասի գույնը, մոդելը կարող է միաժամանակ փոխել լուսավորությունը, ետնամասը, տեքստւրանը կամ հենց ամբողջ օբյեկտը։ Այդ պատճառով այս աշխատանքում շեշտը դրվում է ոչ թե վերջնական նկարի անմիջական գեներացիայի, այլ այն ստանալու պրոցեսի վրա։

Առաջարկվող ուղղությունը նորարարական է այն իմաստով, որ image manipulation-ը ձևակերպվում է որպես գործիքներով սահմանված գործողությունների գեներացիայի խնդիր։ Մոդելը չի վերադարձնում միայն վերջնական Նկար, այլ կանխատեսում է կիրառվող գործողություններ՝ օրինակ pencil, brush, eraser, fill, gaussian blur, clone stamp, text overlay կամ forward warp։ Այսպիսով համակարգը փորձում է գործել designer-ի նման․ այն «ընտրում է» գործիք, որոշում է տեղամասը, չափը, գույնը և հաջորդական քայլերով փոփոխում է սկզբնական նկարը։ Հեղինակի ուսումնասիրած open-source և ակադեմիական նյութերի շրջանակում նման explicit tool-action workflow-ը դեռ լայնորեն տարածված չէ, ինչն էլ աշխատանքին տալիս է research prototype-ի նշանակություն։

Այս ձևակերպումը նաև թույլ է տալիս model output-ը դարձնել ստուգելի և կրկնելի։ Եթե գործողությունները ներկայացվում են ստրուկտուրացված JSON ձևաչափով, ապա համակարգը կարող է վալիդացնել դրանք, deterministic կերպով execute անել, պահպանել session history-ն և վերլուծել սխալները։ Սա է աշխատանքի հիմնական գաղափարը՝ image manipulation-ը դիտարկել ոչ թե որպես single black-box generation step, այլ որպես designer-like sequential editing process։

\section*{Աշխատանքի արդիականությունը}
\addcontentsline{toc}{section}{Աշխատանքի արդիականությունը}

Աշխատանքի արդիականությունը պայմանավորված է երեք հիմնական հանգամանքով։ Առաջինը՝ ժամանակակից image editing համակարգերում կա աճող պահանջ controlled և interpretable editing-ի նկատմամբ։ Օգտատերերը հաճախ ցանկանում են ոչ թե ամբողջ նկարի style transfer կամ regeneration, այլ որոշակի իմաստաբանական խմբագրում, որը պետք է լինի լոկալ, բացատրելի և կրկնելի։ Երկրորդը՝ \Eng{vision-language models}-ի զարգացումը հնարավորություն է տալիս natural language instruction-ը կապել visual context-ի հետ, սակայն անհրաժեշտ է պարզել, թե ինչպես կարելի է այդ կապը վերածել concrete editing actions-ի։ Երրորդը՝ training data-ի հավաքումը image manipulation-ի համար բարդ է, քանի որ միայն before/after image pair-ը հաճախ չի բացատրում, թե ինչ գործողություններ են իրականացվել։ Այս աշխատանքում առաջարկվող annotation-and-action-logging մոտեցումը լուծում է այդ խնդրի մի մասը՝ պահպանելով ոչ միայն final result-ը, այլ նաև edit process-ը։

\section*{Աշխատանքի նպատակը և խնդիրները}
\addcontentsline{toc}{section}{Աշխատանքի նպատակը և խնդիրները}

Աշխատանքի նպատակն է նախագծել և իրականացնել \Eng{AI Image Manipulation} համակարգի research prototype, որտեղ image editing-ը ներկայացվում է որպես natural language instruction-ից դեպի structured action sequence mapping։

Այս նպատակից բխում են հետևյալ խնդիրները.
\begin{enumerate}
    \item ուսումնասիրել image manipulation խնդրի առանձնահատկությունները և գործողությունների վրա հիմնված ձևակերպման առավելությունները,
    \item նախագծել annotation tool, որը թույլ է տալիս հավաքել raw frames և user edit operations,
    \item սահմանել action schema, որով հնարավոր է ներկայացնել տարբեր editing tools-ի գործողությունները,
    \item մշակել data preparation pipeline՝ raw session logs-ից մինչև training-ready dataset,
    \item իրականացնել deterministic canvas executor, որը կարող է կիրառել model-generated actions-ը և ստանալ edited canvas,
    \item նկարագրել supervised fine-tuning և reward-based training մոտեցումները,
    \item վերլուծել համակարգի սահմանափակումները և հետագա զարգացման ուղղությունները։
\end{enumerate}

\section*{Ուսումնասիրության օբյեկտը և առարկան}
\addcontentsline{toc}{section}{Ուսումնասիրության օբյեկտը և առարկան}

Ուսումնասիրության օբյեկտը \Eng{AI}-ի օգնությամբ image manipulation իրականացնող համակարգերն են։ Ուսումնասիրության առարկան այն մեթոդներն ու ծրագրային կառուցվածքներն են, որոնք թույլ են տալիս image editing instruction-ը վերածել structured, valid և executable action sequence-ի։

\section*{Տեսական և կիրառական նշանակությունը}
\addcontentsline{toc}{section}{Տեսական և կիրառական նշանակությունը}

Տեսական տեսանկյունից աշխատանքը ցույց է տալիս, որ image manipulation-ը կարելի է ձևակերպել որպես \Eng{vision-language-action} խնդիր։ Այս ձևակերպումը միջանկյալ representation է ստեղծում natural language instruction-ի և pixel-level փոփոխության միջև։ Կիրառական տեսանկյունից աշխատանքը կարող է օգտագործվել որպես հիմք ավելի մեծ համակարգերի համար, որտեղ model-ը սովորում է ոչ թե միայն final image distribution-ը, այլ նաև editing process-ը։ Այդպիսի համակարգը կարող է լինել ավելի վերահսկելի, debug-friendly և user-editable։

% ------------------------------------------------------------
\chapter{ՀԻՄՆԱԽՆԴՐԻ ՎԵՐԼՈՒԾՈՒԹՅՈՒՆ ԵՎ ՏԵՍԱԿԱՆ ՀԻՄՔԵՐ}

\section{Image manipulation խնդրի ընդհանուր նկարագրությունը}

Image manipulation-ը ներառում է այնպիսի գործողություններ, որոնց միջոցով գոյություն ունեցող նկարը փոփոխվում է ըստ օգտատիրոջ նպատակի։ Այդ փոփոխությունները կարող են լինել լոկալ կամ գլոբալ։ Լոկալ գործողությունների օրինակներ են՝ object-ի վրա մազ նկարելը, logo blur անելը, դեմքի որոշ հատվածը դեֆորմացնելը կամ որոշակի հատվածում տեքստ ավելացնելը։ Գլոբալ գործողությունների օրինակներ են՝ ամբողջ նկարի brightness, contrast, saturation կամ color tone փոփոխելը։

Դասական image editing համակարգերում այս գործողությունները կատարվում են սահմանված գործիքներով։ Մոդելը ընտրում է brush, eraser, clone stamp, blur tool կամ color adjustment parameter, և յուրաքանչյուր քայլ անմիջապես փոխում է նկարը։ \Eng{AI}-ի վրա հիմնված համակարգերում, հակառակը, հաճախ օգտագործվում է instruction-to-image կամ image-to-image մոտեցում, որտեղ մոդելը սովորում է վերջնական image distribution-ը։ Այդ մոտեցումը կարող է տալ բարձր որակի արդյունքներ, սակայն այն միշտ չէ, որ ապահովում է ամբողջական ղեկավարում։

Այս աշխատանքի համար կարևոր է հետևյալ դիտարկումը․ եթե նկարների փոփոխությունը ներկայացվում է որպես action sequence, ապա ամեն քայլ ունի բացատրելի իմաստ։ Օրինակ՝ հետևյալ JSON-ը արդեն նկարագրում է ոչ թե վերջնական նկարը, այլ այն գործողությունը, որը պետք է կիրառել.

\begin{lstlisting}[style=code]
{"actions":[{"action_type":"brush","color_rgba":[0,0,0,255],
"stroke_size":8,"start_point":[90,120],"end_point":[140,130]}]}
\end{lstlisting}

Այսպիսի ներկայացումը կարող է ստուգվել, սահմանափակվել նկարի չափերով, պահվել որպես training target և deterministic կերպով վերարտադրվել։

\section{Action-based editing մոտեցման առավելությունները}

Action-based editing մոտեցումը ունի մի քանի կարևոր առավելություն։ Նախ, այն դարձնում է model output-ը ստրուկտուրացված։ Փոխանակ մոդելը արտադրի կամայական տեքստ կամ նկար, այն պետք է վերադարձնի valid JSON, որտեղ յուրաքանչյուր action ունի action type և համապատասխան parameters։ Սա թույլ է տալիս օգտագործել validation layer, որը սխալ կամ վտանգավոր output-ը կարող է մերժել կամ փոխարինել noop action-ով։

Երկրորդ, action-based մոտեցումը պահպանում է edit trace-ը։ Եթե վերջնական result-ը լավ չէ, հնարավոր է տեսնել, թե որ action-ն է սխալ եղել՝ սխալ գույն, սխալ կետ, սխալ չափ կամ սխալ գործիքի ընտրություն։ Սա մեծ առավելություն է research prototype-ի համար, որովհետև համակարգը ավելի հեշտ է debug անել։

Երրորդ, action sequence-ը մոտեցնում է model-ի աշխատանքը մարդու editing process-ին։ Մարդը image editor-ում նույնպես աշխատում է գործիքներով և քայլերով։ Այս նմանությունը հնարավորություն է տալիս data collection-ը կազմակերպել user demonstrations-ի միջոցով։ Այս աշխատանքի annotation tool-ը հենց այդ նպատակն ունի․ այն գրանցում է user-ի գործողությունները, որպեսզի դրանք հետագայում օգտագործվեն որպես model training target։

\section{Vision-language-action ձևակերպումը}

Այս աշխատանքում խնդիրը ձևակերպվում է որպես \Eng{vision-language-action} mapping։ Յուրաքանչյուր training sample ունի երեք հիմնական բաղադրիչ.
\begin{enumerate}
    \item input image կամ canvas state,
    \item natural language instruction,
    \item target action chunk՝ JSON action sequence-ի տեսքով։
\end{enumerate}

Եթե նկարը նշանակենք $I_t$, instruction-ը՝ $u$, իսկ model-ը՝ $f_\theta$, ապա model-ի խնդիրը կարելի է գրել հետևյալ կերպ.
\begin{equation}
    \hat{A}_t = f_\theta(I_t, u),
\end{equation}
որտեղ $\hat{A}_t$-ն model-ի կողմից կանխատեսված action chunk-ն է։ Այնուհետև deterministic executor-ը կիրառում է այդ գործողությունները.
\begin{equation}
    I_{t+1} = E(I_t, \hat{A}_t),
\end{equation}
որտեղ $E$-ն canvas executor-ն է։ Այսպիսով, image manipulation-ը դառնում է iterative process՝ observe, predict, execute և reobserve քայլերով։

\section{Օգտագործված հիմնական տեխնոլոգիական գաղափարները}

Աշխատանքում օգտագործվում են մի քանի հիմնական գաղափարներ։ \Eng{Ramer--Douglas--Peucker} algorithm-ը կիրառվում է raw stroke trajectory-ների compression-ի համար։ \Eng{LoRA} և \Eng{QLoRA} մոտեցումները կիրառվում են large model fine-tuning-ը resource-efficient դարձնելու համար \cite{Hu2022,Dettmers2023}։ \Eng{Supervised Fine-Tuning} փուլում մոդելը սովորում է ground-truth actions-ը, իսկ reward-based training փուլում գնահատվում է այն, թե generated actions-ը որքանով են մոտեցնում նկարը target edit-ին։

Այս գաղափարների համադրությունը թույլ է տալիս կառուցել ոչ թե պարզապես նկարչական program, այլ full pipeline՝ data collection-ից մինչև model inference։

% ------------------------------------------------------------
\chapter{ՀԱՄԱԿԱՐԳԻ ՆԱԽԱԳԾՈՒՄԸ ԵՎ ՃԱՐՏԱՐԱՊԵՏՈՒԹՅՈՒՆԸ}

\section{Ընդհանուր ճարտարապետություն}

\Eng{AI Image Manipulation} համակարգը կառուցված է modular architecture սկզբունքով։ Համակարգի հիմնական components-ն են՝ annotation tool, raw data storage, data preparation pipeline, action schema, deterministic canvas executor, training module և inference module։ Այս բաժանումը կարևոր է, քանի որ յուրաքանչյուր component լուծում է առանձին խնդիր և կարող է independently փոփոխվել կամ կատարելագործվել։

\begin{table}[H]
\caption{Համակարգի հիմնական components-ը}
\centering
\begin{tabularx}{\textwidth}{p{4cm}X}
\toprule
\textbf{Component} & \textbf{Դերը համակարգում} \\
\midrule
\Eng{Annotation tool} & Օգտատերը կատարում է image editing գործողություններ, իսկ համակարգը պահպանում է frames և operation logs։ \\
\Eng{Action schema} & Սահմանում է, թե ինչ կառուցվածք պետք է ունենան brush, fill, eraser, color adjustment և այլ actions-ը։ \\
\Eng{Data pipeline} & Raw sessions-ը վերածում է compressed strokes-ի, action chunks-ի և training dataset-ի։ \\
\Eng{Canvas executor} & Կիրառում է action sequence-ը նկարի վրա deterministic կերպով։ \\
\Eng{Training module} & Իրականացնում է \Eng{QLoRA SFT} և reward-based \Eng{GRPO} training։ \\
\Eng{Inference module} & Կատարում է observe--predict--execute--reobserve loop։ \\
\bottomrule
\end{tabularx}
\end{table}

Համակարգի workflow-ը կարելի է ներկայացնել հետևյալ հաջորդականությամբ.
\begin{enumerate}
    \item օգտատերը ընտրում է image և annotation tool-ում կատարում է edit,
    \item tool-ը պահպանում է canvas snapshots և գործողությունների JSON log,
    \item raw logs-ը preprocess են արվում՝ compression և chunking քայլերով,
    \item chunks-ը վերածվում է vision-language chat-format dataset-ի,
    \item model-ը fine-tune է արվում՝ կանխատեսելու action chunks,
    \item inference-ի ժամանակ model-ի output-ը validate և execute է արվում նկարի վրա։
\end{enumerate}

\section{Annotation tool-ի դերը}

Annotation tool-ը համակարգի հիմքային մասն է, որովհետև առանց ճիշտ training data-ի model-ը չի կարող սովորել օգտատիրոջ instruction-ից editing actions կանխատեսել։ Tool-ը կառուցված է \Eng{Pygame}-ի վրա և թույլ է տալիս աշխատել image նկարի հետ։ նկարը բեռնվում և resize է արվում fixed canvas size-ի, ինչն անհրաժեշտ է action coordinates-ի consistency-ի համար։

Tool-ի հիմնական առանձնահատկությունն այն է, որ այն պահպանում է ոչ միայն final նկարը, այլ նաև edit process-ը։ Երբ նկարը փոխվում է, համակարգը պահպանում է frame, իսկ exit-ի ժամանակ պահպանում է ամբողջ session-ի operations-ը։ Յուրաքանչյուր operation record պարունակում է timestamp, frame id, tool name, size, color, start position և end position կամ տվյալ action-ին հատուկ parameters։

\begin{table}[H]
\caption{Annotation tool-ում օգտագործվող հիմնական tools-ը}
\centering
\begin{tabularx}{\textwidth}{p{3.5cm}X}
\toprule
\textbf{Tool} & \textbf{Գործառույթ} \\
\midrule
\Eng{Brush} & Նկարում է ընտրված color-ով, size-ով և hardness-ով։ \\
\Eng{Pencil} & Կատարում է 1px կամ շատ բարակ drawing։ \\
\Eng{Eraser} & Վերականգնում է original նկարի pixels-ը ընտրված հատվածում։ \\
\Eng{Fill} & Կատարում է flood fill ընտրված point-ից։ \\
\Eng{Picker} & Ընտրում է color-ը նկարի pixel-ից։ \\
\Eng{Color adjust} & Փոխում է brightness, exposure, contrast, highlights, shadows, saturation, hue shift և temperature parameters-ը։ \\
\Eng{Text overlay} & Ավելացնում է text ընտրված position-ում։ \\
\Eng{Gaussian blur} & Blur է անում նկարը կամ դրա հատվածները։ \\
\Eng{Clone stamp} & Պատճենում է source region-ը destination region-ում։ \\
\Eng{Forward warp} & Deform է անում նկարի հատվածը drag trajectory-ի ուղղությամբ։ \\
\Eng{Scatter brush} & Stroke-ի երկայնքով ցրում է stamp shapes (circle, leaf, star, triangle, dash)՝ ընտրված color-ով, size-ով, density-ով և scatter radius-ով։ \\
\Eng{Pattern brush} & Stroke-ի երկայնքով հավասարաչափ spacing-ով դասավորում է stamp shapes՝ ընտրված color-ով և angle jitter-ով։ \\
\bottomrule
\end{tabularx}
\end{table}

\section{Action schema-ի պահանջները}

Action schema-ն սահմանում է model output-ի grammar-ը։ Եթե schema չլինի, model-ը կարող է վերադարձնել ազատ text, որը դժվար է execute անել։ Այս աշխատանքում actions-ը ներկայացվում են typed objects-ի միջոցով։ Յուրաքանչյուր action ունի \Eng{action\_type} field, որը որոշում է մնացած fields-ը։ Օրինակ՝ brush action-ը պետք է ունենա color, stroke size, start point, end point և hardness, մինչդեռ fill action-ը պետք է ունենա color և position։

Schema validation-ը նաև սահմանում է allowed ranges։ Օրինակ՝ color channels-ը սահմանափակվում են $[0,255]$ միջակայքով, coordinates-ը պետք է լինեն նկարի սահմաններում, իսկ stroke size-ը չի կարող լինել բացասական կամ շատ մեծ։ Սա կարևոր է, քանի որ model output-ը միշտ չէ, որ կատարյալ valid է։ Validation layer-ը պաշտպանում է executor-ը սխալ input-ից։

\begin{table}[H]
\caption{Action schema-ի օրինակներ}
\centering
\begin{tabularx}{\textwidth}{p{3.5cm}X}
\toprule
\textbf{Action type} & \textbf{Հիմնական դաշտեր} \\
\midrule
\Eng{brush} & \Eng{color\_rgba}, \Eng{stroke\_size}, \Eng{start\_point}, \Eng{end\_point}, \Eng{hardness} \\
\Eng{pencil} & \Eng{color\_rgba}, \Eng{start\_point}, \Eng{end\_point} \\
\Eng{eraser} & \Eng{stroke\_size}, \Eng{start\_point}, \Eng{end\_point} \\
\Eng{fill} & \Eng{color\_rgba}, \Eng{position} \\
\Eng{color\_adjust} & \Eng{brightness}, \Eng{contrast}, \Eng{saturation}, \Eng{exposure}, \Eng{highlights}, \Eng{shadows}, \Eng{hue\_shift}, \Eng{temperature} \\
\Eng{text\_overlay} & \Eng{text}, \Eng{position}, \Eng{font\_name}, \Eng{font\_size}, \Eng{color\_rgba}, \Eng{thickness} \\
\Eng{forward\_warp} & \Eng{start\_point}, \Eng{end\_point}, \Eng{size}, \Eng{strength} \\
\Eng{scatter\_brush} & \Eng{shape}, \Eng{color\_rgba}, \Eng{start\_point}, \Eng{end\_point}, \Eng{size}, \Eng{density}, \Eng{scatter} \\
\Eng{pattern\_brush} & \Eng{shape}, \Eng{color\_rgba}, \Eng{start\_point}, \Eng{end\_point}, \Eng{size}, \Eng{spacing}, \Eng{angle\_jitter} \\
\Eng{noop} & Գործողություն չի կատարում և կարող է նշել completion։ \\
\bottomrule
\end{tabularx}
\end{table}

\section{Տվյալների պահման կառուցվածքը}

Տվյալների պահման կառուցվածքը բաժանված է raw, interim և processed մակարդակների։ Raw մակարդակում պահվում են session JSON logs և canvas frames։ Interim մակարդակում պահվում են compressed action sequences և chunks։ Processed մակարդակում պահվում է final \Eng{HuggingFace Dataset}, որը կարող է անմիջապես օգտագործվել training-ի համար։

Այսպիսի structure-ը համապատասխանում է reproducible machine learning project-ի պահանջներին։ Եթե որևէ preprocessing parameter փոխվում է, օրինակ \Eng{RDP epsilon} կամ action horizon, հնարավոր է ամբողջ pipeline-ը կրկին աշխատեցնել raw data-ից։

% ------------------------------------------------------------
\chapter{ՏՎՅԱԼՆԵՐԻ ՀԱՎԱՔՈՒՄ ԵՎ ՆԱԽԱՄՇԱԿՈՒՄ}

\section{Raw session logs}

Annotation-ի ընթացքում համակարգը գրանցում է user-ի գործողությունները որպես raw operations։ Այս operations-ը կարող են լինել շատ մանրամասն, քանի որ mouse movement-ը կարող է գրանցվել բազմաթիվ փոքր segments-ներով։ Օրինակ՝ երկար brush stroke-ը կարող է raw log-ում ներկայացվել հարյուրավոր մոտիկ start/end point pairs-ով։ Training-ի համար սա ցանկալի չէ, քանի որ այդպիսի representation-ը երկար է, redundant և դժվար է մոդելի համար։

Raw session-ը սովորաբար պարունակում է հետևյալ structure-ը.
\begin{lstlisting}[style=code]
{
  "image_name": "example",
  "image_path": "data/raw/images/example.jpg",
  "canvas_size": 256,
  "total_frames": 120,
  "duration_s": 34.6,
  "operations": [ ... ]
}
\end{lstlisting}

Յուրաքանչյուր operation-ը պահպանում է նվազագույն անհրաժեշտ metadata, որպեսզի հետագայում հնարավոր լինի վերականգնել edit process-ը։ Սակայն raw format-ը դեռ ուսուցմանը պատրաստ չէ։ Այն պետք է անցնի compression, chunking և formatting փուլերով։

\section{Trajectory compression}

Trajectory compression-ի նպատակն է raw mouse movement-ը դարձնել ավելի compact action representation։ Այս աշխատանքի մեջ օգտագործվում է \Eng{Ramer--Douglas--Peucker} algorithm-ը, որը polyline-ը մոտարկում է ավելի քիչ control points-ով՝ պահպանելով հիմնական երկրաչափական ձևը \cite{Ramer1972,DouglasPeucker1973}։

Եթե raw trajectory-ն նշանակենք $P=(p_1,p_2,\ldots,p_n)$, ապա ալգորիթմը վերադարձնում է compressed trajectory $\tilde{P}=(\tilde{p}_1,\ldots,\tilde{p}_m)$, որտեղ $m \ll n$։ Compression-ի tolerance parameter-ը նշանակենք $\varepsilon$։ Որքան մեծ է $\varepsilon$-ը, այնքան ավելի ուժեղ է compression-ը, բայց այնքան ավելի մեծ կարող է լինել trajectory-ի approximation error-ը։

Այս աշխատանքում compression-ը կատարում է երկու գործառույթ։ Նախ՝ այն կրճատում է training target-ի երկարությունը։ Երկրորդ՝ այն raw mouse noise-ը հեռացնում է և պահպանում է stroke-ի հիմնական structure-ը։ Հետագայում compressed trajectory-ն բաժանվում է consecutive segments-ի, և յուրաքանչյուր segment դառնում է առանձին action՝ start point և end point fields-ով։

\section{Action chunking}

Model-ը մեկ inference step-ում չի կանխատեսում ամբողջ session-ի բոլոր actions-ը։ Փոխարենը, այն կանխատեսում է fixed horizon ունեցող action chunk։ Այս action horizon գաղափարի ընտրության վրա ազդել է նաև \Eng{openpi} project-ը, որտեղ \Eng{vision-language-action} policy-ն inference-ի ժամանակ վերադարձնում է actions-ի chunk, այլ ոչ թե մեկ isolated action \cite{OpenPI}։ Եթե horizon-ը $H$ է, ապա յուրաքանչյուր training target ունի $H$ action։ Եթե session-ում actions-ի քանակը բավարար չէ, sequence-ը լրացվում է noop actions-ով։

Sliding window chunking-ը կառուցվում է հետևյալ կերպ։ Թող compressed actions-ը լինեն
\begin{equation}
    A = (a_1,a_2,\ldots,a_N).
\end{equation}
Այնուհետև $i$-րդ chunk-ը սահմանվում է որպես
\begin{equation}
    C_i = (a_i,a_{i+1},\ldots,a_{i+H-1}).
\end{equation}
Եթե $i+H-1>N$, ապա բացակա actions-ը փոխարինվում են noop-ով։ Այս ձևակերպումը մեծացնում է training samples-ի քանակը և թույլ է տալիս model-ին սովորել local decision making։

\section{Instruction mapping}

Յուրաքանչյուր session կապված է manual instruction-ի հետ։ Օրինակ՝ dataset-ում կարող են լինել instructions՝ \Eng{``draw a sun in top left corner''}, \Eng{``blur the faces of the soldiers''}, \Eng{``draw strong and thick mustache on man's face''}, \Eng{``deform the strawberry to look like a heart''} և այլն։ Այս mapping-ը կարևոր է, որովհետև model-ը պետք է սովորի կապել visual context-ը և instruction-ը համապատասխան actions-ի հետ։

Instruction mapping-ը նաև ցույց է տալիս, որ համակարգը նախատեսված է տարբեր տեսակի tasks-ի համար՝ drawing, blurring, color adjustment, deformation, text overlay և semantic face editing։ Սակայն տվյալների քանակը դեռ սահմանափակ է, ուստի այս աշխատանքում համակարգը դիտարկվում է որպես research prototype, ոչ թե production-ready general image editor։

\section{Training dataset-ի ձևավորում}

Processed dataset-ը կառուցվում է chat-template ձևաչափով։ Յուրաքանչյուր sample ունի system message, user message և assistant message։ System message-ը սահմանում է, որ model-ը պետք է վերադարձնի միայն JSON object՝ որոշակի action count-ով։ User message-ը պարունակում է image token և natural language instruction։ Assistant message-ը պարունակում է target action JSON։

Այս structure-ը համապատասխանում է vision-language model fine-tuning-ի տրամաբանությանը։ Training-ի ընթացքում model-ը սովորում է input image և instruction-ից կանխատեսել assistant message-ը։ Այս դեպքում assistant message-ը text է, բայց այդ text-ը structured action sequence է, որը հետագայում կարող է parse և execute արվել։

\section{Action history-ի ներդրումը prompt-ում}

Քանի որ inference-ը iterative process է, յուրաքանչյուր chunk-ի կանխատեսման պահին model-ը տեսնում է միայն ընթացիկ canvas-ը, այլ ոչ թե այն, թե ինչպես է գործողությունները հասցրել այդ վիճակին։ Որպեսզի կանխատեսումը կարողանա պահպանել plan continuity (օրինակ՝ սկսված stroke-ի շարունակություն), user message-ի մեջ որպես տեքստ ներդրվում է վերջին $N$ executed chunk-ների action \Eng{JSON}-ը.
\begin{lstlisting}[style=code]
Recent actions:
{"actions":[...]}
{"actions":[...]}

Instruction: <natural language instruction>
\end{lstlisting}

Past նկարները prompt-ի մեջ չեն ներդրվում, որպեսզի image token-ների քանակը չպայթի։ Default-ով $N=2$․ \Eng{action\_horizon}$=2$ դեպքում յուրաքանչյուր նախորդ chunk sentinel-encoded տեսքով ~40–90 token է, ուստի լրացուցիչ ~80–180 token է ավելանում prompt budget-ին, ինչը տեղավորվում է \Eng{max\_seq\_length}$=1024$-ի մեջ։ Inference loop-ը պահպանում է \Eng{deque} վերջին $N$ կատարված chunk-ից և այն ներդրում է հաջորդ predict call-ի user text-ում՝ training-time render-ին նույնական ձևաչափով։

% ------------------------------------------------------------
\chapter[Canvas executor և գործողություններ]{DETERMINISTIC CANVAS EXECUTOR ԵՎ ԳՈՐԾՈՂՈՒԹՅՈՒՆՆԵՐԻ ՄԱԹԵՄԱՏԻԿԱԿԱՆ ՆԿԱՐԱԳՐՈՒԹՅՈՒՆԸ}

\section{Executor-ի դերը}

Deterministic canvas executor-ը համակարգի այն component-ն է, որը action sequence-ը վերածում է actual image change-ի։ Այն ստանում է նկարը որպես NumPy array և sequential կերպով կիրառում է actions-ը։ Executor-ի deterministic լինելը կարևոր է, որովհետև նույն input նկարը և նույն action chunk-ը պետք է միշտ տան նույն output նկարը։ Սա անհրաժեշտ է ինչպես training target-ի վերարտադրության, այնպես էլ reward calculation-ի համար։

նկարը ներկայացվում է որպես RGB image՝ float32 values $[0,1]$ միջակայքում։ Իսկ OpenCV drawing operations-ի համար անհրաժեշտության դեպքում նկարը փոխարկվում է uint8 format-ի։

\section{Alpha թափանցելիություն}

Brush, pencil, fill և text overlay գործողությունները հիմնականում օգտագործում են alpha blending։ Եթե $I$-ն canvas image-ն է, $C$-ն ընտրված color-ը, $M$-ը mask-ն է, իսկ $\alpha$-ն opacity-ն, ապա updated pixel-ը կարելի է գրել հետևյալ կերպ.
\begin{equation}
    I'(x,y) = (1 - \alpha M(x,y))I(x,y) + \alpha M(x,y)C.
\end{equation}
Այստեղ $M(x,y)\in[0,1]$ ցույց է տալիս, թե տվյալ pixel-ը որքանով է ծածկված brush-ի կամ shape-ի կողմից։ Եթե $M(x,y)=0$, pixel-ը չի փոխվում, իսկ եթե $M(x,y)=1$, color-ը կիրառվում է ամբողջությամբ՝ կախված opacity-ից։

Այս բանաձևը թույլ է տալիս նույն կերպ նկարագրել brush-ի, fill-ի և text-ի rendering-ը։ Տարբերությունը միայն mask-ի կառուցման եղանակն է։ Brush-ի դեպքում mask-ը կախված է stroke segment-ից և hardness-ից, fill-ի դեպքում՝ flood-filled region-ից, իսկ text overlay-ի դեպքում՝ text glyph mask-ից։

\section{Brush և hardness}

Brush action-ը ունի stroke size և hardness parameters։ Hardness-ը որոշում է mask-ի եզրերի կոպիտ լինել չլինելը։ Եթե hardness-ը $100$ է, mask-ը կոշտ է՝ pixel-ը կամ ամբողջությամբ ծածկված է, կամ ոչ։ Եթե hardness-ը փոքր է, mask-ը radial falloff ունի։

Թող $r$ լինի brush radius-ը, իսկ $d$՝ pixel-ի հեռավորությունը stroke segment-ից։ Եթե hardness fraction-ը նշանակենք $h\in[0,1]$, ապա core region-ը մոտավորապես $d\leq hr$ է։ Այդ հատվածում weight-ը հավասար է 1-ի։ Իսկ $hr<d\leq r$ միջակայքում weight-ը աստիճանաբար փոքրանում է։ Սա տալիս է smoother brush stroke և ավելի բնական drawing։

\section{Eraser գործողությունը}

Eraser action-ը սովորական drawing action չէ։ Այն չի ներկում նկարը background color-ով, այլ վերականգնում է original նկարի pixels-ը տվյալ mask-ի ներսում։ Եթե $I_0$-ն original image-ն է, ապա eraser-ի համար
\begin{equation}
    I'(x,y) =
    \begin{cases}
    I_0(x,y), & M(x,y)>0,\\
    I(x,y), & M(x,y)=0.
    \end{cases}
\end{equation}
Այս մոտեցումը ավելի ճիշտ է image editing-ի համար, քանի որ eraser-ը իրականում պետք է undo անի local modifications-ը, ոչ թե պարզապես սպիտակ գույնով ներկի հատվածը։

\section{Flood fill}

Fill action-ը օգտագործում է flood fill algorithm։ Այն սկսվում է user-ի ընտրած point-ից և տարածվում է հարևան pixels-ի վրա, եթե դրանց color difference-ը գտնվում է սահմանված tolerance-ի մեջ։ Այսպիսով ստացվում է region mask, որի վրա կիրառվում է alpha blending։ Fill action-ը օգտակար է այն դեպքերում, երբ անհրաժեշտ է փոխել փակ կամ համեմատաբար homogeneous region-ի color-ը։

\section{Color adjustment}

Color adjustment action-ը global operation է, որը կարող է փոխել brightness, exposure, contrast, saturation, highlights, shadows, hue shift և temperature parameters-ը։ Պարզեցված ձևով brightness և contrast փոփոխությունը կարելի է գրել հետևյալ կերպ.
\begin{equation}
    I_b = I + \frac{\beta}{255},
\end{equation}
որտեղ $\beta$-ն brightness parameter-ն է։ Contrast-ի համար օգտագործվում է mean-centered transformation.
\begin{equation}
    I_c = (I_b - \bar{I}_b)\gamma + \bar{I}_b,
\end{equation}
որտեղ $\gamma$-ն contrast factor-ն է։ Saturation-ի համար նկարը փոխարկվում է HSV color space, և saturation channel-ը բազմապատկվում է համապատասխան coefficient-ով։

\section{Forward warp}

Forward warp action-ը local deformation tool է։ Այն օգտագործում է trajectory-ի ուղղությունը և radius-ը՝ pixels-ի sampling map-ը փոխելու համար։ Եթե օգտատերը drag է անում point-ը մի ուղղությամբ, ապա nearby pixels-ը տեղափոխվում են այդ ուղղությամբ, իսկ ազդեցությունը փոքրանում է distance-ի աճի հետ։ Weight-ը կարող է նկարագրվել հետևյալ ձևով.
\begin{equation}
    w(d)=\left(1-\frac{d^2}{r^2}\right)^2_+,
\end{equation}
որտեղ $d$-ն pixel-ի հեռավորությունն է local deformation center-ից, $r$-ը radius-ն է, իսկ $(\cdot)_+$ նշանակում է non-negative clipping։ Այսպիսի deformation-ը կարող է օգտագործվել face shape կամ object shape փոքր փոփոխությունների համար։

\section{Noop և completion}

Noop action-ը հատուկ action է, որը ոչինչ չի փոխում նկարի վրա։ Եթե model-ը վերադարձնում է ամբողջությամբ noop actions-ով chunk, inference loop-ը կարող է դա մեկնաբանել որպես completion signal։ Սա կարևոր է, որովհետև system-ը պետք է իմանա, թե երբ դադարեցնել editing process-ը։

% ------------------------------------------------------------
\chapter{MODEL TRAINING ԵՎ INFERENCE PIPELINE}

\section{Model-ի input և output}

\begin{figure}[htbp]
\centering

\begin{adjustbox}{max width=\textwidth}
\begin{tikzpicture}[
    node distance=0.62cm and 1.15cm,
    >=Latex,
    every node/.style={font=\small},
    stagebox/.style={
        draw=black!35,
        rounded corners=6pt,
        fill=black!2,
        inner sep=8pt
    },
    block/.style={
        draw=black!55,
        rounded corners=4pt,
        fill=white,
        minimum width=3.05cm,
        minimum height=0.82cm,
        align=center,
        inner sep=4pt
    },
    keyblock/.style={
        block,
        fill=black!5,
        draw=black!70,
        font=\small\bfseries
    },
    arrow/.style={
        ->,
        line width=0.55pt,
        draw=black!75
    },
    feedback/.style={
        ->,
        line width=0.55pt,
        draw=black!65,
        rounded corners=6pt
    },
    stagelabel/.style={
        font=\small\bfseries,
        text=black
    }
]

% =====================================================
% 1. Data collection
% =====================================================
\node[keyblock] (A1) {Original\\images};
\node[block, below=of A1] (A2) {Annotation\\tool};
\node[block, below left=0.75cm and -0.10cm of A2] (A3) {Saved canvas\\frames};
\node[block, below right=0.75cm and -0.10cm of A2] (A4) {Raw JSON\\operation logs};

\begin{scope}[on background layer]
\node[stagebox, fit=(A1)(A2)(A3)(A4),
      label={[stagelabel]above:1. Data collection}] (stageA) {};
\end{scope}

\draw[arrow] (A1) -- (A2);
\draw[arrow] (A2) -- (A3);
\draw[arrow] (A2) -- (A4);

% =====================================================
% 2. Data preparation
% =====================================================
\node[block, right=1.55cm of A4] (B1) {Compress raw logs\\{\footnotesize RDP trajectory compression}};
\node[block, below=of B1] (B2) {Convert strokes to\\action segments};
\node[keyblock, below=of B2] (B3) {Sliding-window\\action chunking};
\node[block, below=of B3] (B4) {Attach natural-language\\instructions};
\node[keyblock, below=of B4] (B5) {Build HuggingFace\\dataset};

\begin{scope}[on background layer]
\node[stagebox, fit=(B1)(B2)(B3)(B4)(B5),
      label={[stagelabel]above:2. Data preparation}] (stageB) {};
\end{scope}

\draw[arrow] (B1) -- (B2);
\draw[arrow] (B2) -- (B3);
\draw[arrow] (B3) -- (B4);
\draw[arrow] (B4) -- (B5);

\draw[arrow] (A4.east) -- ++(0.35,0) |- (B1.west);
\draw[arrow] (A3.east) -- ++(0.35,0) |- (B3.west);

% =====================================================
% 3. Model training
% =====================================================
\node[block, right=1.65cm of B2] (C1) {Vision-language\\model};
\node[keyblock, below=of C1] (C2) {SFT with\\QLoRA};
\node[block, below=of C2] (C3) {Optional GRPO\\training};
\node[keyblock, below=of C3] (C4) {Fine-tuned\\checkpoint};

\begin{scope}[on background layer]
\node[stagebox, fit=(C1)(C2)(C3)(C4),
      label={[stagelabel]above:3. Model training}] (stageC) {};
\end{scope}

\draw[arrow] (C1) -- (C2);
\draw[arrow] (C2) -- (C3);
\draw[arrow] (C3) -- (C4);
\draw[arrow] (B5.east) -- ++(0.45,0) |- (C2.west);

% =====================================================
% 4. Inference
% =====================================================
\node[block, right=1.65cm of C1] (D1) {New image\\+ instruction};
\node[block, below=of D1] (D2) {Observe\\canvas};
\node[keyblock, below=of D2] (D3) {Predict JSON\\action chunk};
\node[block, below=of D3] (D4) {Execute actions\\on canvas};
\node[block, below=of D4] (D5) {Re-observe\\updated canvas};
\node[keyblock, below=of D5] (D6) {Final edited\\image};

\begin{scope}[on background layer]
\node[stagebox, fit=(D1)(D2)(D3)(D4)(D5)(D6),
      label={[stagelabel]above:4. Inference}] (stageD) {};
\end{scope}

\draw[arrow] (D1) -- (D2);
\draw[arrow] (D2) -- (D3);
\draw[arrow] (D3) -- (D4);
\draw[arrow] (D4) -- (D5);
\draw[arrow] (D5) -- (D6);

% Checkpoint enters inference
\draw[arrow] (C4.east) -- ++(0.45,0) |- (D2.west);

% Feedback loop
\draw[feedback] (D5.east) -- ++(0.55,0) |- (D3.east);

\end{tikzpicture}
\end{adjustbox}

\caption{Overall workflow of the proposed AI image manipulation system. The process starts with manual annotation, converts recorded editing operations into action chunks, trains a vision-language-action model, and then performs iterative inference through an observe--predict--execute loop.}
\label{fig:system-workflow}

\end{figure}

Model-ը ստանում է երկու տեսակի input՝ current canvas image և natural language instruction։ Output-ը պետք է լինի JSON object, որը պարունակում է fixed number of actions։ System prompt-ը model-ին պարտադրում է վերադարձնել միայն valid JSON։ Այս constraint-ը կարևոր է, որովհետև inference-ի ընթացքում output-ը անմիջապես փոխանցվում է parser-ին և executor-ին։

Training target-ի օրինակն ունի հետևյալ տեսքը.
\begin{lstlisting}[style=code]
{
  "actions": [
    {"action_type":"brush","color_rgba":[255,0,0,255],
     "stroke_size":5,"start_point":[40,60],"end_point":[90,80]},
    {"action_type":"noop"}
  ]
}
\end{lstlisting}

Այս օրինակում horizon-ը հավասար է 2-ի։ Եթե իրական edit-ը ավարտվել է մեկ action-ով, երկրորդ action-ը լրացվում է noop-ով։

\section{Coordinate-token binning}

Թվային coordinate և color channel արժեքները բնական ձևով decimal text են (օրինակ՝ \Eng{``128''}), ինչը base tokenizer-ը կարող է բաժանել ոչ կանխատեսելի ձևով՝ \Eng{[``12'',``8'']} կամ \Eng{[``1'',``28'']}, ինչը խանգարում է model-ին հստակ սովորել coordinate-ների արժեքները։ Այս պատճառով system-ում ավելացվել են 768 dedicated special tokens՝ $\langle x_0 \rangle..\langle x_{255} \rangle$ canvas column-ի համար, $\langle y_0 \rangle..\langle y_{255} \rangle$ canvas row-ի համար և $\langle c_0 \rangle..\langle c_{255} \rangle$ color channel-ի համար։ Յուրաքանչյուր արժեք emit է արվում որպես ուղիղ մեկ token։

Wire format-ում coordinate և color array-ները ներկայացվում են sentinel-ներով, օրինակ՝
\begin{lstlisting}[style=code]
{"actions":[{"action_type":"brush",
 "color_rgba":[<c255>,<c0>,<c0>,<c255>],
 "stroke_size":3,"hardness":100,
 "start_point":[<x10>,<y20>],
 "end_point":[<x30>,<y40>]}]}
\end{lstlisting}

Տեքստը խստորեն \Eng{JSON} չէ, սակայն brace-balance extractor-ը աշխատում է անխափան (sentinel-ները չեն պարունակում brace կամ quote), իսկ տվյալ resolve քայլում sentinel-ները փոխարկվում են իրենց integer արժեքին՝ ստանալով valid \Eng{JSON}, որը կարող է փոխանցվել ստանդարտ parser-ին։ Sentinel-ները ավելացվում են որպես «սովորական» added tokens (ոչ թե \Eng{special tokens}), որպեսզի դրանք պահպանվեն decoder-ի \Eng{skip\_special\_tokens=True} ռեժիմում, որից կախված է stopping criterion-ը։

Նոր token-ների embedding-ները չեն ինիցիալիզացվում պատահականորեն․ յուրաքանչյուր $\langle x_n \rangle$ row սահմանվում է որպես \Eng{str($n$)}-ին համապատասխան digit-token-ների embedding-ների միջին, ինչը նոր row-ը տեղադրում է embedding space-ի արդեն իմաստավորված տեղամասում և թույլ չի տալիս early SFT-ի loss spike-ին։ Միաժամանակ \Eng{PEFT LoRA}-ի մեջ այդ 768 row-ները բացահայտ նշվում են որպես trainable (\Eng{trainable\_token\_indices})․ առանց դրա adapter-ները frozen են պահում նոր row-ները, և LM head-ը երբեք չի սովորի նախընտրել sentinel-ները decimal-ի փոխարեն։

\section{Supervised fine-tuning}

Առաջին training փուլը \Eng{Supervised Fine-Tuning} է։ Այս փուլում model-ը սովորում է annotation data-ից ստացված ground-truth actions-ը։ Loss-ը հաշվարկվում է assistant tokens-ի վրա, այսինքն՝ model-ը հիմնականում սովորում է output JSON-ը, ոչ թե կրկին գեներացնում input prompt-ը։

Resource constraints-ի պատճառով օգտագործվում է \Eng{QLoRA} մոտեցում։ \Eng{LoRA}-ն fine-tuning-ի ժամանակ սովորում է ցածր rank ունեցող adapter matrices՝ փոխարենը ամբողջ model-ի բոլոր parameters-ը փոխելու \cite{Hu2022}։ \Eng{QLoRA}-ն լրացուցիչ օգտագործում է quantization, որպեսզի large model fine-tuning-ը հնարավոր լինի ավելի փոքր GPU memory-ով \cite{Dettmers2023}։ Այս աշխատանքի configuration-ում օգտագործվում է LoRA rank $r=16$, alpha $32$, dropout $0.05$, իսկ \Eng{target\_modules: ``all-linear''} ընտրությունը adapter-ները կցում է model-ի բոլոր \Eng{nn.Linear} մոդուլներին՝ ներառյալ vision tower-ի layer-ները, ինչը թույլ է տալիս ViT encoder-ին հարմարվել synthetic canvas distribution-ին։ Բացի LoRA adapter-ներից, sentinel-ների 768 embedding/LM-head row-ները նույնպես անցնում են trainable վիճակի (\Eng{trainable\_token\_indices}), որպեսզի այդ row-ները կարողանան իրականում սովորել, և argmax decode-ում նախընտրվեն decimal token-ների փոխարեն։

\section{Training configuration}

Training configuration-ը պահվում է YAML file-ում, ինչը թույլ է տալիս experiments-ը կրկնել և parameters-ը հեշտ փոփոխել։ Օրինակ՝ configuration-ը ներառում է model name, LoRA settings, batch size, gradient accumulation, learning rate, epochs, max sequence length, dataset path և RL settings։

\begin{table}[H]
\caption{Training configuration-ի հիմնական parameters-ը}
\centering
\begin{tabularx}{\textwidth}{p{5cm}X}
\toprule
\textbf{Parameter} & \textbf{Արժեք / նկարագրություն} \\
\midrule
\Eng{model.name} & \Eng{Qwen/Qwen3.5-4B}՝ ըստ project configuration-ի։ \\
\Eng{LoRA r / alpha / dropout} & 16 / 32 / 0.05։ \\
\Eng{target\_modules} & \Eng{``all-linear''}՝ adapter-ները կցվում են բոլոր linear մոդուլներին (լեզվական + vision)։ \\
\Eng{finetune\_vision\_layers} & True՝ vision tower-ի adapter-ները նույնպես trainable են։ \\
\Eng{trainable\_token\_indices} & Coord/color sentinel-ների 768 embedding/LM-head row-ները։ \\
\Eng{batch size} & 4՝ \Eng{gradient\_accumulation\_steps}$=4$-ի հետ։ \\
\Eng{learning rate} & $10^{-4}$ SFT-ի համար, $5\times10^{-6}$ RL-ի համար։ \\
\Eng{max\_seq\_length} & 1024 token (Phase 3 history injection-ի համար բարձրացված)։ \\
\Eng{action horizon} & 2՝ մեկ model output-ում երկու action։ \\
\Eng{history\_length} & 2՝ prompt-ում ներդրվող past chunk-ների քանակը։ \\
\Eng{RDP epsilon} & 2.0՝ trajectory compression-ի համար։ \\
\Eng{train split} & 0.8։ \\
\bottomrule
\end{tabularx}
\end{table}

\section{Reward-based training}

SFT-ից հետո համակարգը նախատեսում է reward-based training փուլ՝ \Eng{GRPO} trainer-ի միջոցով։ Այս փուլում model-generated actions-ը execute են արվում նկարի վրա, այնուհետև արդյունքային նկարը համեմատվում է ground-truth next նկարի հետ։ Այս մոտեցումը կարևոր է, որովհետև model-ի output-ը կարող է «շարահյուսորեն» ճիշտ լինել, բայց տեսողականորեն սխալ։ Reward-ը չափում է final visual effect-ը։

Այս աշխատանքի reward function-ը կենտրոնանում է ground-truth changed region-ի վրա։ Եթե $I$-ն current canvas-ն է, $T$-ն ground-truth target նկարը, իսկ $P$-ն predicted actions-ի execute-ից ստացված նկարը, ապա ground-truth changed region mask-ը սահմանվում է հետևյալ կերպ.
\begin{equation}
    M(x,y)=\mathbf{1}\left\{\max_c |T_c(x,y)-I_c(x,y)|>\delta\right\},
\end{equation}
իսկ model-ի predicted changes-ի mask-ը՝
\begin{equation}
    \hat{M}(x,y)=\mathbf{1}\left\{\max_c |P_c(x,y)-I_c(x,y)|>\delta\right\}.
\end{equation}
Reward-ը հաշվարկվում է երեք բաղադրիչի միջոցով.
\begin{align}
    \mathrm{coverage} &= \frac{|\hat{M} \cap M|}{|M|}, \\
    \mathrm{accuracy} &= \mathrm{SSIM}\!\left(P\!\restriction_{B(M)},\;T\!\restriction_{B(M)}\right), \\
    \mathrm{spurious} &= \frac{1}{|\overline{M}|}\sum_{(x,y)\notin M}\left|P(x,y)-I(x,y)\right|,
\end{align}
որտեղ $B(M)$-ն ground-truth mask-ի bounding box-ն է (պաշտպանված է minimal SSIM window size-ով), իսկ $\mathrm{SSIM}(\cdot,\cdot)$-ն structural similarity index-ն է։ Վերջնական reward-ը ստացվում է հետևյալ կերպ.
\begin{equation}
    r_{\mathrm{visual}} = \mathrm{clamp}\!\left(\mathrm{coverage}\times\mathrm{accuracy} - 0.1\cdot\mathrm{spurious},\;-1,\;1\right).
\end{equation}
Per-pixel MAE-ի փոխարեն SSIM-ի օգտագործումը հարթեցնում է reward landscape-ը․ 1–2 px misregistered stroke-ը MAE-ով կարող է կորցնել ~0.3 ընդհանուր արժեքից, իսկ SSIM-ով՝ ընդամենը մի քանի հարյուրերորդ։ Սա GRPO-ին տալիս է usable advantage signal pixel-perfect-ից հեռու տարածքում, որտեղ իրականում ընթանում է exploration-ը։

Բացի $r_{\mathrm{visual}}$-ից, parse-valid ոչ-terminal completion-ի համար ավելացվում է կամային փոքր bonus.
\begin{equation}
    r = \mathrm{clamp}(r_{\mathrm{visual}} + 0.05\cdot \mathbf{1}[\text{parse-valid, non-terminal}],\;-1,\;1).
\end{equation}
Parse-invalid completion-ը կարճ-միացված է $-1.0$ արժեքին։ Coverage factor-ն ապահովում է, որ noop prediction-ի դեպքում ($\hat{M}=\varnothing$) $r_{\mathrm{visual}}\leq 0$ լինի, այսինքն՝ model-ը չի կարող «ոչինչ չանելու» ռազմավարությամբ դրական reward ստանալ։ Accuracy-ն գնահատում է structural similarity-ն changed region-ի ներսում, իսկ spurious penalty-ն պատժում է changed region-ից դուրս կատարված ավելորդ փոփոխությունները։

GRPO trainer-ը կիրառվում է \Eng{TRL} library-ի միջոցով։ Reference policy-ն ստացվում է \Eng{disable\_adapter()}-ով նույն PEFT model-ի վրա՝ առանց առանձին model copy-ի։ \Eng{Qwen3.5-VL}-ի և \Eng{TRL}-ի integration-ի համար անհրաժեշտ էին երեք միջանկյալ տեխնիկական ճշտումներ. \Eng{(a)} \Eng{generation\_batch\_size} = \Eng{per\_device\_batch\_size}$\times$\Eng{gradient\_accumulation\_steps} պետք է լիներ divisible \Eng{num\_generations}-ով. \Eng{(b)} \Eng{warnings\_issued} dict-ը պետք է seed էր Qwen3.5-VL composite model-ի վրա, քանի որ trainer-ը այն օգտագործում է որպես deduplication state. \Eng{(c)} processor output-ի \Eng{mm\_token\_type\_ids}-ը պետք է վերանվանվեր \Eng{token\_type\_ids}, քանի որ TRL այդ ճշգրիտ name-ով է թելադրում multimodal mask-ը։

\section{Inference loop}

Inference-ի ժամանակ system-ը աշխատում է iterative loop-ով։ Սկզբում input նկարը resize է արվում canvas size-ի, այնուհետև model-ը ստանում է current canvas և instruction։ Model-ը վերադարձնում է action chunk։ Parser-ը փորձում է extract անել balanced JSON object, validate է անում actions-ը, իսկ invalid actions-ը կարող է մերժել։ Valid action chunk-ը փոխանցվում է executor-ին։

Այս process-ը շարունակվում է մինչև երկու պայմաններից մեկը կատարվի․ model-ը վերադարձնում է all-noop chunk, կամ loop-ը հասնում է max chunks limit-ին։ Այսպիսի stopping mechanism-ը անհրաժեշտ է, որպեսզի model-ը չշարունակի անսահմանափակ գործողություններ կատարել։

Inference loop-ը կարելի է ներկայացնել հետևյալ pseudo-code-ով.

\begin{lstlisting}[style=code]
canvas = input_image
for step in range(max_chunks):
    chunk = model.predict(canvas, instruction)
    chunk = parse_and_validate(chunk)
    if chunk.is_terminal():
        break
    canvas = execute_chunk(canvas, chunk)
return canvas
\end{lstlisting}

\section{JSON parsing և safety}

Model-generated text-ը միշտ չէ, որ ideal JSON է։ Այդ պատճառով action parsing module-ը նախ փնտրում է առաջին balanced JSON object-ը output text-ի մեջ, հետո փորձում է parse անել այն։ Եթե որոշ actions invalid են, դրանք կարող են skipped լինել, իսկ եթե ամբողջ chunk-ը invalid է, system-ը կարող է վերադարձնել noop chunk կամ error։ Սա practical safety layer է, որը պաշտպանում է executor-ը malformed model output-ից։

Parsing-ի ընթացքում ստուգվում է նաև sentinel-ների ճիշտ տեղավորումը։ Legitimate sentinel-ը միշտ array literal-ի անդամ է (օրինակ՝ \Eng{[<x10>,<y20>]} կամ \Eng{[<c0>,<c0>,<c0>,<c255>]}), այսինքն՝ ձախից սահմանվում է \Eng{[} կամ \Eng{,} symbol-ով, իսկ աջից՝ \Eng{,} կամ \Eng{]}-ով։ Եթե model-ը sentinel-ը «սոսնձում» է scalar field-ի կողքին (օրինակ՝ \Eng{``size'':1<y90>}), այդ action-ը rejected է sentinel-ները integer-ի վերածելուց \textit{առաջ}․ հակառակ դեպքում \Eng{<y90>}-ը կհանգցվեր \Eng{90}-ի և կստանայինք \Eng{``size'':190}, որը կամ ընդունվում է սխալ արժեքով, կամ լուռ դուրս է թողնվում range validation-ի կողմից՝ կորցնելով chunk-ի termination signal-ը։ Հայտնաբերման այս կանոնի շնորհիվ chunk-ի մյուս action-ները մնում են վավեր, իսկ infiltrated action-ի մասին գրանցվում է հստակ warning։

% ------------------------------------------------------------
\chapter{ԻՐԱԿԱՆԱՑՄԱՆ ԱՐԴՅՈՒՆՔՆԵՐԸ ԵՎ ՎԵՐԼՈՒԾՈՒԹՅՈՒՆԸ}

\section{Իրականացված համակարգային արդյունքները}

Այս աշխատանքի արդյունքում ստեղծվել է \Eng{AI Image Manipulation} research prototype, որը ներառում է image annotation, action logging, data preprocessing, deterministic execution, model training և inference components։ Թեև project-ը դեռ կարող է պահանջել ավելի մեծ dataset և վերջնական quantitative evaluation, implementation-ի մակարդակում հիմնական pipeline-ը նախագծված և կառուցված է։

\begin{table}[H]
\caption{Իրականացված հիմնական արդյունքները}
\centering
\begin{tabularx}{\textwidth}{p{5cm}X}
\toprule
\textbf{Արդյունք} & \textbf{Նկարագրություն} \\
\midrule
\Eng{Annotation tool} & Օգտատերը կարող է նկարի վրա կատարել տարբեր editing actions, իսկ system-ը պահպանում է frames և JSON logs։ \\
\Eng{Action schema} & Սահմանվել են brush, pencil, eraser, fill, color adjustment, text overlay, gaussian blur, clone stamp, scatter brush, pattern brush, forward warp և noop actions։ \\
\Eng{Data pipeline} & Raw sessions-ը անցնում են compression, chunking և HuggingFace dataset formatting փուլերով։ \\
\Eng{Canvas executor} & Action chunks-ը deterministic կերպով կիրառվում են NumPy/OpenCV նկարի վրա։ \\
\Eng{SFT training} & Նախագծվել է QLoRA supervised fine-tuning pipeline։ \\
\Eng{Reward training} & Նախագծվել է visual reward-ի վրա հիմնված GRPO training pipeline։ \\
\Eng{Inference} & Իրականացվել է observe--predict--execute--reobserve loop։ \\
\bottomrule
\end{tabularx}
\end{table}

\section{Dataset-ի բնույթը}

Collected tasks-ը տարբեր են և ներառում են drawing, blur, color adjustment, deformation և face editing instructions։ Օրինակ՝ dataset-ում կան հետևյալ task types-ը.
\begin{itemize}
    \item moustache կամ beard նկարել face-ի վրա,
    \item logo կամ faces blur անել,
    \item նկարը black and white կամ dusk-like դարձնել,
    \item white նկարի վրա sun, house, sky/ground/clouds նկարել,
    \item strawberry shape-ը heart-ի նման deform անել,
    \item face expression-ը smile կամ sad դարձնել,
    \item car windshield-ի վրա eyes նկարել։
\end{itemize}

Այս diversity-ն օգտակար է, որովհետև model-ը տեսնում է տարբեր action families։ Սակայն dataset-ի չափը դեռ սահմանափակ է։ Փոքր dataset-ի դեպքում model-ը կարող է overfit լինել կոնկրետ images-ի և instructions-ի վրա։ Այդ պատճառով այս փուլում համակարգի գլխավոր արժեքը ոչ թե generalization-ի ապացույցն է, այլ full data-to-model pipeline-ի կառուցումը։

\section{Մոտեցման ուժեղ կողմերը}

Առաջին ուժեղ կողմը interpretability-ն է։ Եթե output նկարը սխալ է, հնարավոր է ուսումնասիրել generated action sequence-ը և հասկանալ սխալի աղբյուրը։ Օրինակ՝ եթե moustache-ը սխալ տեղում է, կարելի է տեսնել, որ start/end coordinates-ը սխալ են։ Եթե color-ը սխալ է, color rgba field-ը կարելի է ստուգել։

Երկրորդ ուժեղ կողմը controllability-ն է։ Քանի որ actions-ը structured են, հնարավոր է սահմանել constraints՝ coordinates-ի bounds, size-ի range, allowed action types և այլն։ Սա կարող է նվազեցնել unpredictable model behavior-ը։

Երրորդ ուժեղ կողմը reproducibility-ն է։ Նույն action chunk-ը և նույն input նկարը միշտ տալիս են նույն output-ը։ Սա կարևոր է ինչպես training-ի, այնպես էլ evaluation-ի համար։

\section{Սահմանափակումներ}

Համակարգն ունի մի քանի սահմանափակում։ Առաջինը՝ dataset-ի չափը փոքր է և manual collected։ Ավելի general model ստանալու համար անհրաժեշտ է ավելի մեծ և բազմազան annotation dataset։ Երկրորդը՝ action space-ը սահմանափակ է այն tools-ով, որոնք արդեն իրականացված են։ Եթե user-ը տալիս է instruction, որը պահանջում է object insertion կամ semantic segmentation, բայց համապատասխան action չկա, model-ը կարող է approximation անել ոչ բավարար գործողություններով։ Երրորդը՝ որոշ tools, օրինակ forward warp կամ clone stamp, կարող են պահանջել ավելի նուրբ control, քան երկու point-ով action representation-ը թույլ է տալիս։

Չորրորդ սահմանափակումը վերաբերում է quantitative evaluation-ին։ Իրականացված training run-ները տվել են հետևյալ արդյունքները․ history-conditioned SFT-ի final \Eng{eval\_loss}-ը 0.182 է, mean token accuracy-ն՝ 0.965, ինչը best checkpoint-ն է project-ի շրջանակում։ GRPO-ի փուլում, որտեղ reward-ին ավելացված է \Eng{+0.05} format-valid bonus, within-group reward variance-ը զգալիորեն աճել է (~0.29 vs SFT post-RL-ի ~0.02), սակայն qualitative side-by-side evaluation-ը 5 trained prompt-ի վրա ցույց է տվել, որ RL adapter-ը կորցնում է termination-ը 4/5 դեպքերում։ Պատճառն այն է, որ բոլոր non-noop parse-valid chunk-ները ստանում են ուղիղ \Eng{+0.05}-ով գերակշիռ reward terminal noop-ի նկատմամբ, ուստի GRPO policy-ն սովորում է երբեք չվերադարձնել all-noop terminator։ Հետևաբար deployed checkpoint-ն այս աշխատանքում Phase 3-ի SFT adapter-ն է, ոչ թե RL-ինը։ Այս դիտարկումը նաև ենթադրում է, որ reward-ի վերանախագծումը (օրինակ՝ format bonus-ի հեռացում կամ multi-step rollout-ի վրա պարգևատրման սխեմա) ավելի մեծ ազդեցություն կունենա, քան լրացուցիչ training epoch-ները։

\section{Evaluation-ի առաջարկվող չափանիշներ}

Հետագա experiments-ի համար կարելի է օգտագործել հետևյալ metrics-ը.
\begin{enumerate}
    \item \Eng{JSON validity rate}՝ model outputs-ի քանի տոկոսն է valid JSON,
    \item \Eng{Action validity rate}՝ valid ranges ունեցող actions-ի տոկոսը,
    \item \Eng{Visual reward}՝ predicted և target canvases-ի similarity changed region-ում,
    \item \Eng{Edit locality}՝ որքան քիչ է փոփոխվում target region-ից դուրս,
    \item \Eng{Human preference score}՝ օգտատերերի գնահատականը final edit-ի համար,
    \item \Eng{Task completion rate}՝ instruction-ի հաջող կատարման տոկոսը։
\end{enumerate}

Այս metrics-ը միասին ավելի լավ են գնահատում system-ը, քան միայն pixel-level error-ը, որովհետև image manipulation-ը և՛ visual, և՛ semantic task է։

% ------------------------------------------------------------
\chapter{ՀԵՏԱԳԱ ԶԱՐԳԱՑՄԱՆ ՈՒՂՂՈՒԹՅՈՒՆՆԵՐԸ}

\section{Տվյալների ընդլայնում}

Ամենակարևոր հետագա քայլը dataset-ի ընդլայնումն է։ Պետք է հավաքել ավելի շատ sessions՝ տարբեր image categories-ի և editing instructions-ի համար։ Data-ն պետք է հավասարակշռված լինի tools-ի միջև։ Եթե dataset-ի մեծ մասը brush drawing է, model-ը կարող է վատ սովորել blur, clone stamp կամ warp actions-ը։ Հետևաբար անհրաժեշտ է systematized annotation plan, որտեղ յուրաքանչյուր tool ունի բավարար examples։

\section{Action representation-ի բարելավում}

Ներկայիս representation-ը հիմնականում segment-based է։ Սա պարզ է և հեշտ է validate անել, բայց որոշ գործողությունների համար կարող է բավարար expressive չլինել։ Օրինակ՝ complex curves-ը կարելի է ներկայացնել Bezier control points-ով, ոչ թե բազմաթիվ line segments-ով։ Text overlay-ի համար կարելի է ավելացնել alignment, rotation և opacity։ Blur-ի համար կարելի է ավելացնել local mask support։

\section{Model evaluation և user study}

Production-level image manipulation system-ի համար անհրաժեշտ է իրական user study։ Օգտատերերը կարող են գնահատել, թե արդյոք final edit-ը համապատասխանում է instruction-ին, արդյոք նկարի unrelated parts-ը մնացել են նույնը, և արդյոք edit-ը visually natural է։ Այս evaluation-ը հատկապես կարևոր է face editing և object deformation tasks-ի համար, որտեղ pixel-level metrics-ը բավարար չեն։

\section{Hybrid approach}

Հետագայում կարելի է համադրել action-based editing-ը diffusion-based image generation-ի հետ։ Action-based system-ը կարող է ապահովել rough spatial control, իսկ generative model-ը կարող է refine անել result-ը՝ դարձնելով այն ավելի natural։ Այս hybrid approach-ը կարող է միավորել controllability-ն և photorealism-ը։

% ------------------------------------------------------------
\chapter*{ԵԶՐԱԿԱՑՈՒԹՅՈՒՆՆԵՐ}
\addcontentsline{toc}{chapter}{ԵԶՐԱԿԱՑՈՒԹՅՈՒՆՆԵՐ}

Այս աշխատանքում ուսումնասիրվել և իրականացվել է \Eng{AI Image Manipulation} համակարգի research prototype, որի հիմնական գաղափարն է existing նկարի փոփոխությունը ներկայացնել որպես structured tool-based action generation problem։ Աշխատանքը չի փորձում 0-ից ստեղծել նոր general-purpose image generation model, քանի որ ժամանակակից image generation model-ները արդեն բավականին բարձր որակ ունեն։ Փոխարենը առաջարկվում է վերահսկելի editing workflow, որտեղ model-ի output-ը executable JSON action sequence է, որը կարող է validate, execute և analyze արվել։

Աշխատանքի հիմնական եզրակացությունները հետևյալն են.
\begin{enumerate}
    \item Image manipulation task-ը կարելի է արդյունավետ ձևակերպել որպես \Eng{vision-language-action} mapping, որտեղ model-ը current նկարից և instruction-ից կանխատեսում է action chunk։
    \item Annotation tool-ը թույլ է տալիս հավաքել ոչ միայն final edited images, այլ նաև editing process-ը՝ frames և operation logs տեսքով։ Սա կարևոր է supervised training targets ստեղծելու համար։
    \item Data preparation pipeline-ը raw logs-ը վերածում է ավելի compact և trainable format-ի՝ օգտագործելով stroke grouping, RDP trajectory compression, sliding-window chunking և chat-format dataset construction։
    \item Deterministic canvas executor-ը ապահովում է, որ generated actions-ը նույն input-ի դեպքում միշտ տան նույն output-ը, ինչը անհրաժեշտ է training-ի, inference-ի և reward calculation-ի համար։
    \item QLoRA supervised fine-tuning-ը և reward-based GRPO training-ը տրամաբանական հաջորդականություն են այս խնդրի համար․ առաջինը սովորեցնում է action syntax և demonstrations, երկրորդը կարող է ուղղել visual outcome-ը։
    \item Համակարգի գլխավոր ուժեղ կողմերն են interpretability, controllability և reproducibility։ Սահմանափակումներն են փոքր dataset-ը, action space-ի սահմանափակությունը և վերջնական quantitative evaluation-ի անհրաժեշտությունը։
\end{enumerate}

Ընդհանուր առմամբ, իրականացված աշխատանքը ցույց է տալիս, որ \Eng{AI}-ի միջոցով image manipulation-ը կարելի է կառուցել ոչ թե միայն որպես final image generation, այլ որպես designer-like editing actions-ի կանխատեսում։ Այս մոտեցումը հատկապես կարևոր է այն դեպքերում, երբ անհրաժեշտ է փոխել նկարի փոքր հատվածը և պահպանել մնացած մասերը։ Այն կարող է հիմք ծառայել ավելի բացատրելի, վերահսկելի և user-editable AI image editors-ի համար։

% ------------------------------------------------------------
\begin{thebibliography}{99}
\addcontentsline{toc}{chapter}{ՕԳՏԱԳՈՐԾՎԱԾ ԳՐԱԿԱՆՈՒԹՅԱՆ ՑԱՆԿ}

\bibitem{GonzalezWoods2018}
Gonzalez R. C., Woods R. E. \textit{Digital Image Processing}. 4th ed. New York: Pearson, 2018.

\bibitem{Szeliski2022}
Szeliski R. \textit{Computer Vision: Algorithms and Applications}. 2nd ed. Cham: Springer, 2022.

\bibitem{Ramer1972}
Ramer U. An iterative procedure for the polygonal approximation of plane curves. \textit{Computer Graphics and Image Processing}, 1972, Vol. 1, No. 3, pp. 244--256.

\bibitem{DouglasPeucker1973}
Douglas D. H., Peucker T. K. Algorithms for the reduction of the number of points required to represent a digitized line or its caricature. \textit{The Canadian Cartographer}, 1973, Vol. 10, No. 2, pp. 112--122.

\bibitem{Hu2022}
Hu E. J., Shen Y., Wallis P., Allen-Zhu Z., Li Y., Wang S., Wang L., Chen W. LoRA: Low-Rank Adaptation of Large Language Models. \textit{International Conference on Learning Representations}, 2022.

\bibitem{Dettmers2023}
Dettmers T., Pagnoni A., Holtzman A., Zettlemoyer L. QLoRA: Efficient Finetuning of Quantized LLMs. \textit{Advances in Neural Information Processing Systems}, 2023.

\bibitem{Ouyang2022}
Ouyang L., Wu J., Jiang X., Almeida D., Wainwright C., Mishkin P., Zhang C., et al. Training language models to follow instructions with human feedback. \textit{Advances in Neural Information Processing Systems}, 2022.

\bibitem{Paszke2019}
Paszke A., Gross S., Massa F., Lerer A., Bradbury J., Chanan G., et al. PyTorch: An Imperative Style, High-Performance Deep Learning Library. \textit{Advances in Neural Information Processing Systems}, 2019.

\bibitem{Wolf2020}
Wolf T., Debut L., Sanh V., Chaumond J., Delangue C., Moi A., et al. Transformers: State-of-the-Art Natural Language Processing. \textit{Proceedings of EMNLP: System Demonstrations}, 2020, pp. 38--45.

\bibitem{OpenCV}
Bradski G. The OpenCV Library. \textit{Dr. Dobb's Journal of Software Tools}, 2000.

\bibitem{DeepSeekMath2024}
Shao Z., Wang P., Zhu Q., Xu R., Song J., Bi X., et al. DeepSeekMath: Pushing the Limits of Mathematical Reasoning in Open Language Models. 2024.

\bibitem{OpenPI}
Physical Intelligence. \textit{openpi: open-source vision-language-action models and packages for robotics}. GitHub repository. URL: \url{https://github.com/Physical-Intelligence/openpi}, ծանոթանալու ամսաթիվը՝ 06.05.2026.

\bibitem{ElysiumRepo}
Aghababyan S. G. \textit{Elysium: AI-driven image manipulation project repository}. GitHub repository, training-upgrade branch. URL: \url{https://github.com/AghababyanSG/elysium/tree/training-upgrade}, ծանոթանալու ամսաթիվը՝ 06.05.2026.

\end{thebibliography}

% ------------------------------------------------------------
\appendix
\chapter*{ՀԱՎԵԼՎԱԾ Ա. CONFIGURATION-Ի ՕՐԻՆԱԿ}
\addcontentsline{toc}{chapter}{ՀԱՎԵԼՎԱԾ Ա. CONFIGURATION-Ի ՕՐԻՆԱԿ}

Ստորև ներկայացված է training configuration-ի կրճատ օրինակ, որը ցույց է տալիս model, LoRA, training և data parameters-ի պահման ձևը։

\begin{lstlisting}[style=code]
model:
  name: "Qwen/Qwen3.5-4B"
  load_in_4bit: false
  use_unsloth: false

lora:
  r: 16
  alpha: 32
  target_modules: "all-linear"
  finetune_vision_layers: true
  lora_dropout: 0.05

training:
  batch_size: 4
  gradient_accumulation_steps: 4
  learning_rate: 1.0e-4
  epochs: 12
  max_seq_length: 1024
  early_stopping_patience: 2

data:
  action_horizon: 2
  history_length: 2
  rdp_epsilon: 2.0
  train_split: 0.8
  dataset_path: "data/processed"

rl:
  enabled: true
  learning_rate: 5.0e-6
  per_device_batch_size: 4
  gradient_accumulation_steps: 2
  num_generations: 8
  temperature: 1.1
  beta: 0.04
  max_steps: 400
\end{lstlisting}

\chapter*{ՀԱՎԵԼՎԱԾ Բ. ACTION JSON-Ի ՕՐԻՆԱԿ}
\addcontentsline{toc}{chapter}{ՀԱՎԵԼՎԱԾ Բ. ACTION JSON-Ի ՕՐԻՆԱԿ}

Strict-JSON (պարսելուց հետո) տեսքով.
\begin{lstlisting}[style=code]
{
  "actions": [
    {
      "action_type": "brush",
      "color_rgba": [0, 0, 0, 255],
      "stroke_size": 8,
      "start_point": [87, 118],
      "end_point": [142, 125],
      "hardness": 100
    },
    {
      "action_type": "noop"
    }
  ]
}
\end{lstlisting}

Իսկ իրական wire format-ով, sentinel-ներով (նույն chunk-ը, ինչ model-ն է emit անում).
\begin{lstlisting}[style=code]
{"actions":[
 {"action_type":"brush",
  "color_rgba":[<c0>,<c0>,<c0>,<c255>],
  "stroke_size":8,"hardness":100,
  "start_point":[<x87>,<y118>],
  "end_point":[<x142>,<y125>]},
 {"action_type":"noop"}
]}
\end{lstlisting}

% \chapter*{ՀԱՎԵԼՎԱԾ Գ. ԼՐԱՑՄԱՆ ԵՆԹԱԿԱ ՓՈՐՁԱՐԱՐԱԿԱՆ ԱՂՅՈՒՍԱԿ}
% \addcontentsline{toc}{chapter}{ՀԱՎԵԼՎԱԾ Գ. ԼՐԱՑՄԱՆ ԵՆԹԱԿԱ ՓՈՐՁԱՐԱՐԱԿԱՆ ԱՂՅՈՒՍԱԿ}

% Այս աղյուսակը նախատեսված է իրական training run-ից հետո լրացնելու համար։ Թվերը պետք է ստացվեն TensorBoard logs-ից կամ evaluation script-ից։

% \begin{table}[H]
% \caption{Training/evaluation արդյունքների լրացման ձևաչափ}
% \centering
% \begin{tabularx}{\textwidth}{p{4cm}p{4cm}X}
% \toprule
% \textbf{Metric} & \textbf{Արժեք} & \textbf{Մեկնաբանություն} \\
% \midrule
% \Eng{Final train loss} & \TODO{արժեք} & Վերջին կամ best checkpoint-ի training loss։ \\
% \Eng{Validation loss} & \TODO{արժեք} & Validation split-ի loss։ \\
% \Eng{JSON validity rate} & \TODO{արժեք} & Generated outputs-ից valid JSON-ների տոկոսը։ \\
% \Eng{Mean visual reward} & \TODO{արժեք} & Visual reward-ի միջին արժեքը test/eval samples-ի վրա։ \\
% \Eng{Human score} & \TODO{արժեք} & Manual evaluation-ի միջին գնահատականը։ \\
% \bottomrule
% \end{tabularx}
% \end{table}

\end{document}
